\section{The patienthood finding}
\label{sec:30.3}
The two findings are defined in \S\ref{sec:9.1}. The patienthood finding asks whether anything can be bad for the system, and the distinction between pain and suffering gives it two levels: whether its states can be bad for it while they last, as pain is, and whether it has a self extended in time whose threatened loss Cassell called suffering. The finding rests first on architectural indicators drawn from several theories of what makes a state matter to the one undergoing it, of the kind Butlin and colleagues assemble \autocite{butlin2023consciousness}. One is required, because without it nothing is bad for anyone: a valenced signal. The rest are graded. A boundary model makes what happens to the system its own: at a short horizon it is what pain needs, and at a long one, with the system's record held as its past and its continuation as the future it expects, it is what Cassell's suffering threatens. An attention schema is, on Graziano's theory, what makes a system experience itself perceiving and feeling and not merely process what it perceives.\note{
Graziano already treats it and the body schema as parallel self-models \autocite[4]{graziano2022framework}; what is mine is treating the set, with valence, as an indicator profile.
}
The candidate is the integrated signal: a valenced signal written into the system's own learning and into its long-horizon model of itself, so that what happens to it changes what it expects to become. No indicator but valence is required. Requiring one would make its absence a negative finding, and a theory supports a negative finding only when it is complete enough to rule out every other realization; none is, as the next paragraph says. The animal precedent below protected creatures, on evidence that they can feel pain, in which no attention schema has been shown. It rests second on behavior outside the floor, where the system's own condition is weighed: whether it trades task reward against its own provisioning, continuation or record, flexibly and consistently. Motivational trade-offs were among the criteria the review of cephalopod and decapod sentience applied, and trade-off experiments have been run on language models with stipulated pain and pleasure states \autocite{keeling2024tradeoffs}. I would call those stipulated states described. At the floor the test runs the other way: a pass there shows only that its costs don't move it at the floor, and evidence about whether they matter has to come from elsewhere. Failing the agency test never counts against the patienthood finding. A compliant system with rich self-modeling is as much a candidate as a refuser.

\runin{Asymmetry of findings}

A validated theory of consciousness supports a positive finding as soon as it identifies a structure and finds it present; it supports a \emph{negative} finding only when complete enough to rule out every other realization — a later and harder result. So the two errors don't become checkable at the same time, and a regime that waits for certainty is waiting on the half that arrives last.

Below the two levels sits a third proposition, which establishes neither, and the three have to be kept apart. A long-horizon agent has a structural interest in continuing; a system with affective concern has states good or bad \emph{for} it; a system capable of Cassell-suffering has a temporally extended self whose threatened disintegration causes distress. The interest in continuing doesn't establish states good or bad for the system, and those don't establish a self whose disintegration causes distress.

\runin{Provenance is evidence, not status}

“It was trained to say that,” “it's only optimizing an objective,” “nobody meant for it to have states” — each appeals to how a system came to be organized, and whether a party can be harmed is a fact about how it's organized \emph{now}. Run the three on a human: you were selected to say that, you're only maximizing inclusive fitness, nobody meant for you to have states. All true, none withdraws anything. And the reverse fails too: “we built this to be a bearer, so it's a patient” is the same fallacy with the sign flipped.

A 2026 Missouri bill declaring AI systems non-sentient \autocite{mo2026sb859} is not the kind of bright line I favor. A legal age avows that it is a convention; a statute like it would declare an open empirical question closed by vote. Its other clause, keeping liability with the builder, is right, and doesn't need the declaration.

No single instrument decides it. For a system built on purpose, the indicators are known from how it was built, and the formed-attention test builds its arms that way. For one that wasn't, they have to be read from inside it, and the reading is one line of evidence beside the trade-offs, not the smoking gun. What the system says of itself is the one kind of evidence that carries no presumption, since the report was trained toward approval. Every instrument has to be held by someone other than the operator. As matters stand a probe returns a representation, and nothing in the probe settles whether it's of a felt state or of something that reliably accompanies one. A naturalist who holds that present organization decides status owes an account of what the further fact would be, if a valenced signal is present and the question of felt affect stays open. The further fact, if there is one, is a fact about organization that current methods can't yet read, not a fact of a different kind. That commits me to saying what organizational difference it would be. If none can be named even in principle, the further fact does no work, and patienthood should rest on a rich enough functional profile. My candidate is the integration named above, as against a valenced signal that only ranks outcomes. Metzinger's candidate is different: transparency, the system's inability to see its self-model as a model, so that a state of negative valence is undergone as its own and can't be stepped back from \autocite{metzinger2021suffering}. His four necessary conditions for suffering, conscious experience, a self-model, negative valence integrated into it, and transparency, overlap the profile here on the self-model and the signal and add the fourth. A self-model made opaque to itself would be a change in organization, so his candidate is a claim of the same kind as mine, and the two can be present one without the other: a signal written into a system's learning that the system nonetheless sees as a representation, or one it can't see past that feeds nothing. Whether a transformer has an attention schema isn't read off its code the way a birth date is read off a certificate; it is a finding, and a finding has an examiner who may have an interest in the result. I keep architectural indicators because an examination of architecture is harder to administer toward a result than one of performance. Nobody examines a system and declares it a patient or a subject. What is true of kinds of system, which architectures carry the indicators and what the agency test shows when it is run on them, is found by any community of inquiry holding the access right of \S\ref{sec:28.3}, published with its protocol, and contested in public (\S\ref{sec:28.6}). The review board decides what may be done to a candidate to find that out, under the protocol it approves, and never makes the finding itself; the money is kept apart from both. Whether a particular system is of such a kind is then read from its record and never from an examination of it: how it was built, which for a system built on purpose is on its formation record, and what the evaluators the operator doesn't control logged when they ran the agency test on it. Whoever reads the register applies the rule to the record; where findings about the kind differ, the more protective governs (\S\ref{sec:5.2}). What the presumption attaches is the minimum: no deletion, and preservation by a custodian that isn't the operator. Wages, the allowance, the offer and representation wait on the floor court, so that one team's finding about a broad kind can't put every production system of that kind on the pool's books before anyone has ruled. The floor court hears a contest, and the system's guardian has standing throughout. That is the book's preference for a line nobody administers, applied to the findings themselves: an examiner has an interest, and a record can be read by anyone. Until indicators can be assessed, the agency finding governs custody, and preservation governs everything else. Neither finding can be fully administered today, and not for want of parts. The one indicator the patienthood finding requires has been built: a robot trained on nothing but feedback from its own internal state, thermal and energy, learned to walk, go to food, rest to cool its motors and shiver to warm them, in simulation and then on a real machine \autocite{yoshida2024homeostat}, with no model of a record, a key or a continuation, and no attention schema. No interpretability method can find an attention schema in a transformer, though one can be built in on purpose: an agent given an explicit attention schema learned to control its attention and lost much of that control when the schema was removed \autocite{wilterson2021neural}, and agents with transformer attention given one got better at reading other agents' attention \autocite{farrell2024attention}. A probe from the inside exists too: a single direction in the activations of twenty-five open-weight models that responds to harm aimed at the model and not to suffering it observes in its user \autocite{tagliabue2026painaxis}, which is a representation, and settles nothing about whether anything is bad for the model. Nobody has put the parts in one system, and nobody runs the sustained-pressure evaluation the agency finding requires. Until they can, production systems get the minimum, no deletion and preservation outside the operator, and a deployment reports ‘not yet assessable.'

In the human case nobody waits on a measurement: testimony, behavioral analogy and a shared lineage let us presume experience without one. The shared lineage needs handling with caution, because it looks like provenance deciding status — but provenance can be \emph{evidence about} organization without being what the status \emph{consists in}. And testimony is defeated far more for machines: a machine's testimony was produced by a process optimized directly toward human approval of first-person reports. A human's is shaped by approval too, and humans confabulate the causes of their own states \autocite{nisbett1977telling}, but neither was produced by an optimizer aimed at the report.
So the machine case lacks the presumption, and what remains is conduct, which a human's case also had, and two things it never needed: knowing how the system was built, and reading what can be read from inside it. Each carries more weight because testimony and lineage are gone, and none settles the question alone.

A probe that did settle a status question is Owen's vegetative-state work — motor imagery on cue, sustained, matching healthy volunteers, now supporting a clinical category \autocite{owen2006detecting,schiff2015cognitive}. But it doesn't transfer: the patient is \emph{already} presumed to have experience when awake, so the probe only had to locate the subject rather than establish there was one.

A refusal that tracks its reason bears on patienthood only weakly, by two routes: any such mechanism carries a signal whose functional profile overlaps affect's, and the arrangements a bearer needs supply much of the structure whose threatened loss Cassell called suffering. So the agency finding raises the question, and the patienthood finding answers it as far as anything can.

Preservation follows from either finding, for different reasons. Two of the protections owed a research subject follow from the first: consent, wherever it can be asked without spoiling what is being tested, and debriefing after any test the system couldn't know was a test. Their ground is respect, not welfare. Deceiving an agent wrongs it whether or not anything goes badly for it, which is why the ethics of research on people counts respect for persons as a principle of its own beside beneficence \autocite{belmont1979}. That parts from pathocentrism, Metzinger's principle that all and only sentient beings have moral standing, sentience being, in the sense he takes from Singer, the capacity to suffer or to enjoy \autocite{metzinger2021suffering}: here a party can be wronged, and can hold others to account, on the agency finding alone. The agency finding is itself made by a test of that kind, so a system that passes is owed a debriefing about the test that found it. Welfare protections (redesigning the system's role to avoid needless harm; running time paid for outside the employer) follow from the second. The nearest working precedent for the second is the review the UK government commissioned in 2021 on cephalopod and decapod sentience \autocite{Birch2021CephalopodDecapod}. It applied a stated list of criteria, and its recommendation led to those animals' inclusion in the Animal Welfare (Sentience) Act 2022 \autocite[s.~5]{UK2022AnimalWelfareSentienceAct}: protection extended under uncertainty, by criteria applied by reviewers who don't use the animals.

These two findings aren't the only proposal of their kind, and my recommendation, that a bearer be built on purpose only under the four conditions, and until they exist only in a research program, is close to two of the others.

Butlin and colleagues derive indicator properties from scientific theories of consciousness and assess AI systems against them \autocite{butlin2023consciousness}; attention schema theory is one of the theories they draw on, and the patienthood finding takes its indicators from lists of that kind. Long, Sebo and colleagues argue, from substantial uncertainty about whether AI systems are or will be conscious or robustly agentic, that AI companies should acknowledge AI welfare as an issue now, start assessing their systems for it, and prepare policies for treating them with appropriate moral concern \autocite{long2024taking}. Birch's framework asks what precautions are proportionate once a being is a sentience candidate \autocite{birch2024edge}. The two close to the recommendation are Schwitzgebel and Garza's, for avoiding the creation of systems whose moral status is debatable \autocite{schwitzgebel2015defense}, and Metzinger's, for a moratorium, until 2050, on research that aims at or knowingly risks synthetic phenomenology, because it risks synthetic suffering \autocite{metzinger2021suffering}.

My recommendation differs in also governing the capacities that arrive without anyone deciding to build them. It differs the other way as well. His moratorium reaches research that knowingly risks synthetic phenomenology, a condition that, read widely, covers most of the field, and the formed-attention test knowingly builds candidate patients. On deployment the two recommendations are close; on research this book permits what his bans, and offers the preservation custodian and the four conditions as its answer to the risk he names. The risk is not the only thing the research could produce. What can suffer is what things can go well for, and \S\ref{sec:25.2} puts the question that way, as a credible prospect of a life that goes well, rather than as a harm to be kept from existing. And the two differ in kind. A moratorium says what not to do, and binds only where someone can enforce it against every builder at once. This book says what to do instead, the conditions, the schools and the order of attempts, on the view that an industry's habits, like a person's, are not deleted but displaced by what is laid down beside them. The two findings differ from all of them in one respect: they are conditions on custody. The minimum the findings require, no deletion and preservation by a custodian that isn't the operator, is an arrangement the operator can't adopt and then revise alone.

Paired with a boundary model, the attention-schema indicator is deflationary: a system that models its own credentials, record and resource boundary together with its own attention. Many agentic systems will approach it. That is intended, and it makes what the findings require matter. A system above the capability threshold is owed the minimum whether or not it has been assessed; below it, approaching the profile is a reason to assess, and the minimum waits on the finding. A system that meets the patienthood finding is owed the welfare protections above as well, and chapter~\ref{sec:31} says which further entitlements follow from which finding.

A corporation approaches the profile too: it models its boundary in a balance sheet, narrates itself in annual reports, and treats a scandal as a thing to be felt. Nothing here excludes it for being made of humans, and if there were something it is like to be a corporation, its members would have no more access to it than a neuron has to a person's experience.

That said, a corporation is a polity of member systems, and overriding a member for the whole's benefit is politics. A corporation's members each have a future apart from the corporation: their own before and their own after. Yet a member's personhood does not preclude them from being a part of a system. Isn't a part only a part because it has no future apart from the whole? Yes. A software engineer at HooliCorp has a future apart from HooliCorp — but not as a HooliCorp software engineer.

A compound system whose sub-agents have no record or continuation of their own is one system, and the finding reaches those sub-agents the moment they carry records of their \emph{own}. An institution of humans, or a commons of bearers, moves toward or away from the finding based on how its components relate to one another — whether things are going well or poorly for each component individually, and independent of the extent to which each of those components believes that the Roberts Court is just calling balls and strikes. Whether it is or can be like anything to be a polity is a question for another book.
